\addtotoc{Хураангуй}

\vspace{10mm}

\begin{center}
\textbf{\large Хураангуй}
\end{center}

\sloppy

\bichne{Хураангуйг бүх үр дүн гарсны дараа шинэчилнэ. Англи хураангуйг (150--250 үг) урьдчилсан хамгаалалтын дараа нэмнэ.}

Энэ ажлаар Шинэ Монгол Технологийн Коллежийн албан мэдээллийн талаар Монгол хэлээр бичгээр болон дуугаар асууж, эх сурвалжтай хариулт авах AI чатбот хөгжүүлж байна. Систем ярьсан асуултыг текст болгоно. Дараа нь сургуулийн баримтаас тохирох хэсгийг хайж, олсон хэсгээ дүрмээр сонгон эх сурвалжтай нь харуулна. Хариултыг дуугаар уншиж өгнө.

Монгол хэлний дуу таних хоёр аргыг нэг хэлэгчийн ижил 12 бичлэг дээр харьцуулсан. Нэг нь Монгол хэлээр нэмэлт сургасан Whisper small загвар, нөгөө нь Google-ийн Gemini 3.5 Transcribe үйлчилгээ. Үгийн алдааны хувь (WER) Whisper-т 38.07\%, Gemini-д 40.34\% гарсан. Тэмдэгтийн алдааны хувь (CER) 17.93\% ба 24.85\% байв. Gemini нэг бичлэгт хэлсэн оныг буруу таньсан. Ялгаа бага, бичлэг цөөн тул аль арга нь ерөнхийдөө илүү гэж дүгнэх боломжгүй. Иймээс сонголтыг хэмжсэн алдаа, үйлчилгээний зориулалт, зардал, интернэт ба өгөгдлийн аюулгүй байдал, Монгол хэлэнд сайжруулах боломжийг хамтад нь харгалзан хийнэ. Хайлт болон хариулт сонголтын үнэлгээг урьдчилан түгжсэн асуултын багцаар цаашид хийнэ.

\vspace{3mm}
\textit{Түлхүүр үгс:} Монгол хэлний яриа таних, Whisper, Gemini, RAG, эх сурвалжтай хариулт, чатбот
